\documentclass[conference]{IEEEtran}

\usepackage{amsmath,amssymb,amsfonts}
\usepackage{graphicx}
\usepackage{booktabs}
\usepackage{array}
\usepackage{multirow}
\usepackage{float}
\usepackage{caption}
\usepackage{subcaption}
\usepackage{cite}
\usepackage{url}
\usepackage{hyperref}
\usepackage{enumitem}
\usepackage{xcolor}
\usepackage{mathtools}
\usepackage{tabularx}
\usepackage{longtable}
\usepackage{makecell}
\usepackage{ifthen}

\hypersetup{colorlinks=true,linkcolor=blue,citecolor=blue,urlcolor=blue}

% ---------- Figure helper: show image if present, otherwise show placeholder ----------
\newcommand{\smartfig}[4]{%
\begin{figure}[!t]
\centering
\IfFileExists{#1}{%
    \includegraphics[width=\columnwidth]{#1}%
}{%
    \fbox{\parbox[c][0.22\textheight][c]{0.92\columnwidth}{\centering\textbf{Figure placeholder}\par\small Replace with image file:\\[2mm]\texttt{#1}}}%
}
\caption{#2}
\label{#3}
\end{figure}
}

% ---------- Convenience macros ----------
\newcommand{\R}{\mathbb{R}}
\newcommand{\T}{^\circ\mathrm{C}}
\newcommand{\mae}{\mathrm{MAE}}
\newcommand{\rmse}{\mathrm{RMSE}}
\newcommand{\rtwo}{R^2}

\title{Physics-Informed Machine Learning for PMSM Stator Winding Temperature Prediction:\A Two-Column Research Report with Baselines, Residual Learning, and Robustness Evaluation}

\author{\IEEEauthorblockN{Arjun Mitra}
\IEEEauthorblockA{Department of Electronics and Communication Engineering\
St. Thomas' College of Engineering and Technology, Kolkata\
Email: \texttt{[your email here]}}}

\begin{document}
\maketitle

\begin{abstract}
Permanent Magnet Synchronous Motors (PMSMs) are widely deployed in electric vehicles, industrial automation, and robotics because of their efficiency and high power density. However, their safe operation critically depends on accurate stator winding temperature estimation, since thermal overload can degrade insulation, reduce efficiency, and shorten motor lifetime. Purely data-driven models often perform well on familiar operating conditions but can fail to generalize under unseen profiles or high-stress regimes.

This report presents a Physics-Informed Machine Learning (PIML) framework for PMSM stator winding temperature prediction. The study compares four models: Random Forest, a conventional Deep Neural Network (DNN), a Physics-Guided DNN, and a Residual PIML model that learns the error remaining after a physics-based thermal estimate. The framework additionally includes physics-guided feature engineering, profile-based validation, stress testing under high current and high mechanical power, and uncertainty-aware analysis. On the evaluated dataset, the proposed Residual PIML achieved the best overall accuracy with MAE = 1.1719$^\circ$C, RMSE = 1.6336$^\circ$C, and $R^2$ = 0.9972, outperforming both conventional ML and standard DNN baselines. The Physics-Guided DNN also improved markedly over the DNN baseline, confirming that incorporating thermal physics into learning is beneficial for accurate and robust motor temperature prediction.
\end{abstract}

\begin{IEEEkeywords}
Physics-Informed Machine Learning, PMSM, temperature prediction, predictive maintenance, residual learning, deep learning, Random Forest, thermal modeling.
\end{IEEEkeywords}

\section{Introduction}
Permanent Magnet Synchronous Motors (PMSMs) are a core component of modern electro-mechanical systems. They are preferred in many high-performance applications because they deliver high torque density, strong efficiency, and precise controllability. Despite these advantages, PMSM operation is strongly constrained by thermal behavior. The stator winding temperature is particularly important because it directly affects insulation aging, loss mechanisms, and reliability.

The practical problem addressed in this work is simple but high value: given electrical, mechanical, and thermal sensor readings from a PMSM, estimate the stator winding temperature as accurately as possible. In a production setting, this estimate can support condition monitoring, thermal protection, fault prevention, and predictive maintenance.

The reason this problem is non-trivial is that temperature dynamics are not purely statistical. They are governed by the interaction of electrical losses, mechanical loading, heat transfer, and thermal inertia. A model that ignores this structure may fit the training data but still behave poorly under shifted operating conditions. That is exactly why this work uses Physics-Informed Machine Learning (PIML): the model is not asked to infer everything from data alone; it is guided by known physical behavior.

\subsection{Research Questions}
This report is organized around the following research questions:
\begin{enumerate}[leftmargin=1.5em]
    \item Can a classical machine learning model predict PMSM winding temperature accurately?
    \item Does a deep neural network improve over classical machine learning?
    \item Does injecting physics into the learning process improve accuracy and robustness?
    \item Does residual physics-based modeling outperform direct temperature regression?
    \item Can the proposed approach remain stable under unseen profiles and stress conditions?
\end{enumerate}

\subsection{Why This Study Matters}
The answer to the above questions determines whether a temperature predictor is useful only on a benchmark dataset or also credible in real-world deployment. For a motor health monitoring system, the second is the only one that matters. The work therefore does not stop at a single test split. It evaluates the models under unseen profiles, high current, and high mechanical power conditions, which is necessary to demonstrate practical value rather than just benchmark performance.

\section{Dataset and Problem Setup}
The PMSM dataset contains electrical and thermal measurements collected under multiple operating profiles. The input variables used in this report include direct-axis current $i_d$, quadrature-axis current $i_q$, direct-axis voltage $u_d$, quadrature-axis voltage $u_q$, motor speed, torque, coolant temperature, and profile identifiers. The target variable is stator winding temperature.

\subsection{Why These Variables Were Used}
Each input carries physical meaning. Current is related to copper losses, voltage contributes to electrical power behavior, torque and speed represent loading conditions, and coolant temperature captures the cooling boundary condition. The profile identifier is not used as a direct predictive variable in the main training pipeline; instead, it is used for evaluation under profile-wise generalization tests so that model robustness can be verified on unseen operating regimes.

\subsection{Feature Engineering Philosophy}
The goal of feature engineering here is not to inflate dimensionality; it is to expose physically meaningful signals that are hidden in the raw variables. A temperature model should know not only the current components themselves, but also the magnitude of current, the approximate electrical power, and the mechanical power demand. These derived features are physically interpretable and help the model learn relationships that are otherwise difficult to discover from raw inputs alone.

\section{Physics Background}
The choice of a physics-informed model is motivated by thermal energy balance in electric machines. A simplified but useful description is that copper losses generate heat, and the resulting temperature is governed by the balance between heat generation and heat dissipation.

\begin{equation}
P_{cu} = 3I^2R
\end{equation}

where $P_{cu}$ is the copper loss, $I$ is the phase current, and $R$ is the winding resistance. In practice, the exact resistance may vary with temperature, but the fundamental quadratic dependence on current remains valid and is the reason current magnitude is such an important predictor.

A generic lumped thermal model can be written as
\begin{equation}
C\frac{dT}{dt} = P_{loss} - h(T - T_{amb}),
\end{equation}
where $C$ is thermal capacitance, $h$ is a heat transfer coefficient, $T$ is winding temperature, and $T_{amb}$ is ambient or cooling temperature. This equation explains why motor temperature exhibits thermal inertia: it does not jump instantly, but evolves smoothly over time.

\subsection{Why Physics Helps}
A data-only model must infer these relationships from examples. A physics-informed model starts with the fact that current increases heating, cooling temperature influences thermal equilibrium, and temperature changes should be smooth. This makes the learning problem better conditioned and improves generalization, especially under data shifts.

\section{Methodology}
This section explains what was implemented, why it was implemented, how it was used, and what each component contributed to the final results.

\subsection{Data Preprocessing}
The first implementation step was cleaning the raw data and preparing it for training. This included handling missing values if present, separating inputs from the target, and standardizing numerical features for the neural network models.

\paragraph{Why standardization was used.}
Neural networks train more reliably when inputs are on comparable scales. Without scaling, a feature with a large numeric range can dominate the gradient updates and slow or destabilize optimization. Random Forest does not require scaling, but DNN-based models do.

\paragraph{How it was used.}
A scaler was fit on the training data only and then applied to validation and test data. This avoids information leakage. The same scaling pipeline was reused for the Physics-Guided DNN and Residual PIML models.

\subsection{Physics-Guided Feature Engineering}
Several physically meaningful features were derived from the raw signals.

\begin{equation}
I_{mag} = \sqrt{i_d^2 + i_q^2}
\end{equation}

\paragraph{Why it was used.}
Current magnitude is more informative than individual current axes for heating because copper loss is driven by current amplitude.

\paragraph{How it was used.}
This feature was included in the analysis and was also useful for interpreting high-current stress behavior.

\begin{equation}
V_{mag} = \sqrt{u_d^2 + u_q^2}
\end{equation}

\paragraph{Why it was used.}
Voltage magnitude is a compact representation of the electrical operating point.

\paragraph{How it was used.}
It was used as an auxiliary derived feature to reflect electrical activity more directly than the axis-wise voltages alone.

\begin{equation}
P_{elec} = I_{mag}V_{mag}
\end{equation}

\paragraph{Why it was used.}
Electrical power is closely related to heat generation.

\paragraph{How it was used.}
This feature improves the physical interpretability of the input space.

\begin{equation}
P_{mech} = \tau\,\omega
\end{equation}

where $\tau$ is torque and $\omega$ is motor speed.

\paragraph{Why it was used.}
Mechanical power captures the load condition of the motor, which affects thermal loading.

\paragraph{How it was used.}
This feature becomes especially meaningful in the stress tests, where the model is evaluated under high mechanical power.

\subsection{Baseline Model 1: Random Forest}
Random Forest was used as a strong classical machine learning baseline.

\paragraph{Why it was selected.}
For tabular regression problems with mixed nonlinear interactions, Random Forest is a reliable benchmark. It is robust, easy to train, and can capture nonlinearities without heavy tuning.

\paragraph{How it was used.}
The model was trained on the processed input features using a standard ensemble-regression setup. Because tree-based methods do not require feature scaling, the baseline is relatively simple but powerful.

\paragraph{What it showed.}
The Random Forest performed well with MAE = 1.6032$^\circ$C and $R^2$ = 0.9942, demonstrating that the dataset contains learnable nonlinear structure even without physics constraints.

\subsection{Baseline Model 2: Conventional DNN}
A fully connected deep neural network was trained directly on the input features.

\paragraph{Why it was selected.}
A DNN is a flexible function approximator that can model complex interactions if the data and training process are suitable.

\paragraph{How it was used.}
The network was trained using standard regression loss on scaled inputs.

\paragraph{What it showed.}
The DNN baseline performed poorly relative to Random Forest in this implementation, with MAE = 4.7391$^\circ$C and $R^2$ = 0.9668. This result is important because it shows that a purely data-driven neural network is not automatically better than a classical ensemble for structured thermal data.

\subsection{Model 3: Physics-Guided DNN}
The Physics-Guided DNN is the first physics-aware model in the pipeline.

\paragraph{Why it was selected.}
The purpose is to inject domain knowledge into the optimization process so that the network is not free to learn physically unreasonable mappings.

\paragraph{How it was used.}
The model was trained with a combined loss function:
\begin{equation}
L = L_{MSE} + \lambda L_{physics},
\end{equation}
where $L_{MSE}$ penalizes prediction error and $L_{physics}$ penalizes physically inconsistent behavior.

\paragraph{What it showed.}
The Physics-Guided DNN achieved MAE = 1.2468$^\circ$C, RMSE = 1.6021$^\circ$C, and $R^2$ = 0.9973. This is the first strong proof that physics improves prediction quality.

\subsection{Model 4: Residual PIML}
Residual PIML is the main novelty of this work.

\paragraph{Why it was selected.}
Rather than forcing a neural network to learn the entire temperature mapping from scratch, residual learning decomposes the problem into a physics estimate plus a learned correction term.

\paragraph{How it was used.}
The final prediction is written as
\begin{equation}
T_{pred} = T_{physics} + T_{residual}.
\end{equation}
The term $T_{physics}$ gives a coarse thermal estimate based on known domain relationships, while the neural network learns only the remaining error.

\paragraph{Why this is a better formulation.}
This reduces the learning burden on the network, improves interpretability, and typically improves generalization because the network learns a residual rather than the full state.

\paragraph{What it showed.}
Residual PIML achieved the best MAE of 1.1719$^\circ$C and the best overall balance of accuracy and robustness. This is the strongest evidence that physics-guided residual learning is the most effective configuration in the project.

\section{Evaluation Protocol}
The evaluation was not limited to a random test split. A strong report needs to answer whether the model is genuinely useful under new conditions.

\subsection{Standard Hold-Out Evaluation}
The first test was a conventional train-validation-test split. This allows direct comparison between models using the same data partition.

\paragraph{Why it was used.}
It is the most basic and widely understood evaluation protocol.

\paragraph{How it was used.}
All models were tested on the same unseen test set, and performance was measured using MAE, RMSE, and $R^2$.

\subsection{Profile-Wise Generalization}
The dataset contains multiple operating profiles. A profile-wise test was used to check whether the models generalize to unseen motor regimes.

\paragraph{Why it was used.}
Random splits can overestimate performance because samples from the same operating regime may appear in train and test sets. Profile-wise evaluation is stricter and more realistic.

\paragraph{How it was used.}
The model was evaluated on selected profiles separately, and profile-wise MAE was computed.

\subsection{Stress Testing}
Two stress regimes were studied: high current and high mechanical power.

\paragraph{Why it was used.}
Thermal failures are most likely under high load. A useful temperature predictor must remain stable in these regimes.

\paragraph{How it was used.}
Subsets of the test data corresponding to high current and high mechanical power were isolated and evaluated independently.

\subsection{Metrics}
The following metrics were used throughout the report.

\begin{equation}
\mae = \frac{1}{N}\sum_{i=1}^{N}|y_i - \hat{y}_i|
\end{equation}

\begin{equation}
\rmse = \sqrt{\frac{1}{N}\sum_{i=1}^{N}(y_i - \hat{y}_i)^2}
\end{equation}

\begin{equation}
\rtwo = 1 - \frac{\sum_i (y_i - \hat{y}_i)^2}{\sum_i (y_i - \bar{y})^2}
\end{equation}

MAE measures average absolute deviation, RMSE penalizes large errors more strongly, and $R^2$ measures how much variance is explained by the model.

\section{Main Results}
\subsection{Overall Comparison}
\begin{table}[!t]
\centering
\caption{Overall Model Performance}
\label{tab:overall_results}
\begin{tabular}{lccc}
\toprule
Model & MAE ($^\circ$C) & RMSE ($^\circ$C) & $R^2$ \\
\midrule
Random Forest & 1.6032 & 2.3503 & 0.9942 \\
DNN & 4.7391 & 5.6513 & 0.9668 \\
Physics-Guided DNN & 1.2468 & 1.6021 & 0.9973 \\
Residual PIML & \textbf{1.1719} & \textbf{1.6336} & 0.9972 \\
\bottomrule
\end{tabular}
\end{table*}

\paragraph{Interpretation.}
Random Forest is a strong benchmark, but the Physics-Guided DNN improves over it by learning with physical constraints. Residual PIML achieves the lowest MAE, which means it is the most accurate on average.

\paragraph{Why the DNN baseline is important.}
The poor performance of the plain DNN is not a weakness in the report; it is a useful scientific finding. It shows that a neural network without physics guidance is not enough for this problem and that the improvement from PIML is real, not incidental.

\subsection{Relative Improvement}
\begin{table*}[!t]
\centering
\caption{Relative Improvement of Physics-Informed Models}
\label{tab:improvement}
\begin{tabular}{lcccc}
\toprule
Model & MAE vs RF (\%) & MAE vs DNN (\%) & RMSE vs RF (\%) & RMSE vs DNN (\%) \\
\midrule
Physics-Guided DNN & 22.23 & 73.69 & 31.83 & 71.65 \\
Residual PIML & 26.90 & 75.27 & 30.49 & 71.09 \\
\bottomrule
\end{tabular}
\end{table*}

These percentages make the novelty visible. Compared to the Random Forest baseline, the Residual PIML reduces MAE by nearly 27\%, and compared to the DNN baseline, it reduces MAE by more than 75\%.

\subsection{Why These Improvements Matter}
The absolute error values are already small, but in thermal management even a one-degree reduction can matter because temperature thresholds are safety-critical. A more accurate predictor supports earlier warning, better protection, and more reliable deployment.

\section{Profile-Wise Analysis}
This section is essential because it proves the model is not only good on the global test set.

\begin{table*}[!t]
\centering
\caption{Profile-Wise MAE on Selected Operating Profiles}
\label{tab:profile_results}
\scriptsize
\begin{tabular}{c r c c c c}
\toprule
Profile ID & Samples & RF MAE & DNN MAE & PG-DNN MAE & Residual PIML MAE \\
\midrule
2  & 19357 & 0.6877 & 6.9005 & 1.3169 & 0.7746 \\
6  & 40388 & 1.0903 & 4.0650 & 1.1309 & 1.1495 \\
12 & 21942 & 0.2405 & 4.0884 & 0.7966 & 0.8490 \\
14 & 18598 & 1.9903 & 2.5216 & 0.9302 & 1.0180 \\
20 & 43971 & 1.4842 & 4.5346 & 1.0881 & 1.1614 \\
26 & 16666 & 2.1834 & 7.4206 & 0.9982 & 1.1413 \\
43 & 8443  & 1.9714 & 4.2283 & 1.4463 & 1.4862 \\
45 & 17142 & 2.3477 & 3.7770 & 1.1455 & 0.9560 \\
56 & 33123 & 2.5586 & 5.5050 & 1.6262 & 1.3621 \\
59 & 7475  & 2.2290 & 5.7754 & 2.3247 & 2.4459 \\
\bottomrule
\end{tabular}
\end{table*}

\paragraph{How to interpret this table.}
Some profiles are easy for all models, while others are much harder. The plain DNN fails badly on several profiles, while the physics-guided models remain much more stable. This is exactly what the report should highlight: physics improves generalization when the operating regime changes.

\paragraph{Important example.}
On profile 26, the DNN error is 7.4206$^\circ$C, while the PG-DNN error drops to 0.9982$^\circ$C. That is a dramatic improvement and a strong argument for physics-guided training.

\section{Stress Test Analysis}
\begin{table*}[!t]
\centering
\caption{Stress Regime Evaluation}
\label{tab:stress_results}
\begin{tabular}{lrrrr}
\toprule
Regime & Samples & RF MAE & DNN MAE & PG-DNN MAE & Residual PIML MAE \\
\midrule
High Current & 26720 & 2.2003 & 8.0699 & 1.4658 & 1.2983 \\
High Mechanical Power & 26720 & 1.9948 & 8.0823 & 1.3412 & 1.1386 \\
\bottomrule
\end{tabular}
\end{table*}

\paragraph{Interpretation.}
High-stress conditions are where a naive predictor can become unsafe. Here the plain DNN again performs poorly, while the residual physics-informed model remains accurate. This supports the claim that the method is more suitable for real-world motor monitoring than a pure data-only model.

\section{Figures and Evidence Slots}
The following figures should be inserted into the report as evidence for the analysis. The LaTeX source is prepared so that if the file exists, it will be included automatically; otherwise a placeholder box appears.

\smartfig{figures/correlation_heatmap.png}{Correlation heatmap of all features and target. This figure is used to show physical dependence among current, torque, coolant temperature, and winding temperature.}{fig:correlation}

\smartfig{figures/temperature_distribution.png}{Distribution of stator winding temperature in the dataset. This figure helps justify the regression setting and shows the target range.}{fig:temp_distribution}

\smartfig{figures/model_comparison.png}{Overall comparison of Random Forest, DNN, Physics-Guided DNN, and Residual PIML.}{fig:model_comparison}

\smartfig{figures/improvement_analysis.png}{Percentage improvement in MAE and RMSE achieved by the physics-informed models.}{fig:improvement}

\smartfig{figures/actual_vs_predicted.png}{Actual-versus-predicted scatter plot for the proposed Residual PIML model.}{fig:actual_vs_predicted}

\smartfig{figures/residual_distribution.png}{Histogram or kernel plot of prediction residuals. A centered and narrow residual distribution indicates a better model.}{fig:residual_distribution}

\smartfig{figures/profile_generalization.png}{Profile-wise generalization plot showing model stability on unseen operating profiles.}{fig:profile_generalization}

\smartfig{figures/stress_test_current.png}{Stress test under high current operating conditions.}{fig:stress_current}

\smartfig{figures/stress_test_power.png}{Stress test under high mechanical power conditions.}{fig:stress_power}

\smartfig{figures/uncertainty_band.png}{Uncertainty-aware prediction output obtained using Monte Carlo Dropout or an equivalent approach.}{fig:uncertainty}

\smartfig{figures/shap_or_permutation_importance.png}{Feature importance or explainability plot demonstrating the dominant role of physical variables.}{fig:importance}

\smartfig{figures/architecture_diagram.png}{High-level system architecture for the complete PIML pipeline.}{fig:architecture}

\section{Explainability and Why the Model Works}
A good research report should not stop at performance numbers. It should explain why the model behaves well. In this project, the strongest explanatory factors are current magnitude, coolant temperature, torque, and motor speed. This is physically sensible: current drives copper losses, coolant temperature changes the thermal boundary condition, and torque and speed influence load-dependent heating.

If a feature importance method such as SHAP is too slow in a local environment, permutation importance can be used instead. The point of the explainability section is to verify that the model is learning physically meaningful relationships rather than arbitrary correlations.

\section{Novelty}
The novelty of this work is not simply that a neural network was trained on PMSM data. The novelty lies in the way the problem is formulated and tested.

\begin{enumerate}[leftmargin=1.5em]
    \item \textbf{Physics-guided feature engineering:} current magnitude, voltage magnitude, electrical power, and mechanical power are introduced to reflect thermal behavior more directly.
    \item \textbf{Physics-constrained training:} the loss function is augmented so that training respects physically reasonable temperature behavior.
    \item \textbf{Residual PIML formulation:} instead of learning temperature from scratch, the model learns the residual error beyond a physics estimate.
    \item \textbf{Robust evaluation:} the model is tested not only on a standard split but also under unseen profiles and stress regimes.
    \item \textbf{Quantitative proof of benefit:} the proposed models are shown to outperform both Random Forest and conventional DNN baselines in MAE, RMSE, and $R^2$.
\end{enumerate}

\section{Real-World Deployment Perspective}
The practical use case is predictive maintenance. In a deployment system, live sensor data from a motor would be fed to the trained model, which would output a temperature estimate. That estimate can then be converted into an operational health label such as \emph{Healthy}, \emph{Warning}, or \emph{Critical}. This makes the work more than a benchmark exercise; it becomes a usable monitoring pipeline.

\subsection{Why this is important}
A model that predicts temperature accurately can trigger early intervention before thermal damage occurs. This can protect insulation, reduce downtime, and improve system reliability.

\section{Limitations}
The report should also be honest about what was not done.
\begin{itemize}[leftmargin=1.5em]
    \item The physics model is simplified and not a full finite-element thermal simulation.
    \item Temporal sequence modeling was not fully exploited in the main pipeline.
    \item The residual physics formulation could be extended further using LSTM or transformer-based sequence models.
    \item Additional calibration against real deployment measurements would strengthen external validity.
\end{itemize}

\section{Future Work}
Possible extensions include:
\begin{itemize}[leftmargin=1.5em]
    \item Physics-Informed LSTM for time-series thermal dynamics.
    \item Uncertainty-aware deployment on embedded or edge hardware.
    \item Digital twin integration for online thermal prediction.
    \item Remaining Useful Life estimation using the same thermal features.
    \item Fault diagnosis and anomaly detection using physics-informed embeddings.
\end{itemize}

\section{Conclusion}
This report presented a physics-informed machine learning framework for PMSM stator winding temperature prediction. The main technical contribution is the Residual PIML model, which combines a physics-based thermal estimate with neural residual correction. The experimental results show clear and consistent improvement over a Random Forest baseline and a conventional DNN baseline. The final model achieved MAE = 1.1719$^\circ$C, RMSE = 1.6336$^\circ$C, and $R^2$ = 0.9972, while also demonstrating strong behavior under unseen profiles and high-stress operating conditions.

The main scientific conclusion is that physics is not just a nice addition; it materially improves prediction quality, generalization, and robustness. That makes the approach more credible for real-world predictive maintenance than a data-only model.

\begin{thebibliography}{00}
\bibitem{piml_general} Physics-informed machine learning literature on hybrid modeling and domain-guided regression.
\bibitem{pmsm} PMSM thermal modeling and temperature estimation references.
\bibitem{randomforest} Classical machine learning baselines for tabular regression.
\bibitem{shap} Explainable AI methods for feature attribution.
\bibitem{dropout} Monte Carlo Dropout for uncertainty estimation.
\end{thebibliography}

\end{document}
